\chapter{2011年北京卷（文）}

\section{单选题}

\begin{problem}
已知全集 \(U = \R\)，集合 \(P = \{x \mid x^2 \leqslant 1\}\)，那么 \(\complement_U P =\)（\quad）

\choices
    {\((-\infty, -1]\)}
    {\([1, +\infty)\)}
    {\((-1, 1)\)}
    {\((-\infty, -1) \cup (1, +\infty)\)}
\end{problem}

\begin{answer}
D
\end{answer}

\begin{problem}
复数 \(\frac{\i - 2}{1 + 2\i} =\)（\quad）

\choices
    {\(\i\)}
    {\(-\i\)}
    {\(-\frac{4}{5} - \frac{3}{5}\i\)}
    {\(-\frac{4}{5} + \frac{3}{5}\i\)}
\end{problem}

\begin{answer}
A
\end{answer}

\begin{problem}
如果 \(\log_{\frac{1}{2}} x < \log_{\frac{1}{2}} y < 0\)，那么（\quad）

\choices
    {\(y < x < 1\)}
    {\(x < y < 1\)}
    {\(1 < x < y\)}
    {\(1 < y < x\)}
\end{problem}

\begin{answer}
D
\end{answer}

\begin{problem}
若 \(p\) 是真命题，\(q\) 是假命题，则（\quad）

\choices
    {\(p \wedge q\) 是真命题}
    {\(p \vee q\) 是假命题}
    {\(\neg p\) 是真命题}
    {\(\neg q\) 是真命题}
\end{problem}

\begin{answer}
D
\end{answer}

\begin{problem}
某四棱锥的三视图如图所示，该四棱锥的表面积是（\quad）

\bitmapfigure[max width=0.58\linewidth]{img/2011/beijing_liberal/q5_fig1.png}

\choices
    {\(32\)}
    {\(16 + 16\sqrt{2}\)}
    {\(48\)}
    {\(16 + 32\sqrt{2}\)}
\end{problem}

\begin{answer}
B
\end{answer}

\begin{problem}
执行如图所示的程序框图，若输入 \(A\) 的值为 2，则输出的 \(P\) 的值为（\quad）

\texfigure[max width=0.38\linewidth]{img_repaint/2011/beijing_liberal/q6_fig1.tex}

\choices
    {\(2\)}
    {\(3\)}
    {\(4\)}
    {\(5\)}
\end{problem}

\begin{answer}
C
\end{answer}

\begin{problem}
某车间分批生产某种产品，每批的生产准备费用为 \(800\text{ 元}\). 若每批生产 \(x\) 件，则平均仓储时间为 \(\frac{x}{8}\text{ 天}\)，且每件产品每天的仓储费用为 \(1\text{ 元}\). 为使平均到每件产品的生产准备费用与仓储费用之和最小，每批应生产产品（\quad）

\choices
    {\(60\text{ 件}\)}
    {\(80\text{ 件}\)}
    {\(100\text{ 件}\)}
    {\(120\text{ 件}\)}
\end{problem}

\begin{answer}
B
\end{answer}

\begin{problem}
已知点 \(A(0,2)\)，\(B(2,0)\). 若点 \(C\) 在函数 \(y = x^2\) 的图象上，则使得 \(\triangle ABC\) 的面积为 2 的点 \(C\) 的个数为（\quad）

\choices
    {\(4\)}
    {\(3\)}
    {\(2\)}
    {\(1\)}
\end{problem}

\begin{answer}
A
\end{answer}

\section{填空题}

\begin{problem}
在 \(\triangle ABC\) 中，若 \(b = 5\)，\(\angle B = \frac{\pi}{4}\)，\(\sin A = \frac{1}{3}\)，则 \(a =\)\fillinblank{}.
\end{problem}

\begin{answer}
\(\frac{10\sqrt{2}}{3}\)
\end{answer}

\begin{problem}
已知双曲线 \(x^{2}-\frac{y^{2}}{b^{2}}=1\)（\(b>0\)）的一条渐近线的方程为 \(y = 2x\)，则 \(b=\)\fillinblank{}.
\end{problem}

\begin{answer}
\(2\)
\end{answer}

\begin{problem}
已知向量 \(\bs{a} = (\sqrt{3}, 1)\)，\(\bs{b} = (0, -1)\)，\(\bs{c} = (k, \sqrt{3})\). 若 \(\bs{a} - 2\bs{b}\) 与 \(\bs{c}\) 共线，则 \(k =\)\fillinblank{}.
\end{problem}

\begin{answer}
\(1\)
\end{answer}

\begin{problem}
在等比数列 \(\left\{a_{n}\right\}\) 中，\(a_{1}=\frac{1}{2}\)，\(a_{4}=-4\)，则公比 \(q=\)\fillinblank{}；\(a_{1}+a_{2}+\cdots+a_{n}=\)\fillinblank{}.
\end{problem}

\begin{answer}
\(-2\)；\(\frac{1-(-2)^{n}}{6}\)
\end{answer}

\begin{problem}
已知函数 \(f(x) = \begin{cases}
\dfrac{2}{x}, & x \geqslant 2,\\
(x - 1)^3, & x < 2,
\end{cases}\)
若关于 \(x\) 的方程 \(f(x) = k\) 有两个不同的实根，则实数 \(k\) 的取值范围是\fillinblank{}.
\end{problem}

\begin{answer}
\((0,1)\)
\end{answer}

\begin{problem}
设 \(A(0,0)\)，\(B(4,0)\)，\(C(t + 4,3)\)，\(D(t,3) \)（\(t \in \R\)）. 记 \(N(t)\) 为平行四边形 \(ABCD\) 内部（不含边界）的整点的个数. 其中整点是指横，纵坐标都是整数的点，则 \(N(0) =\)\fillinblank{}；\(N(t)\) 的所有可能取值为\fillinblank{}.
\end{problem}

\begin{answer}
\(6\)；\(6\)，\(7\)，\(8\)
\end{answer}

\section{解答题}

\begin{problem}
已知函数 \( f(x) = 4\cos x\sin \left(x + \frac{\pi}{6}\right) - 1 \).
\begin{enumerate}
\item 求 \( f(x) \) 的最小正周期.
\item 求 \( f(x) \) 在区间 \( \left[-\frac{\pi}{6}, \frac{\pi}{4}\right] \) 上的最大值和最小值.
\end{enumerate}
\end{problem}

\begin{solution}
\begin{enumerate}
\item 最小正周期为 \(\pi\).
\item 在区间 \(\left[-\frac{\pi}{6},\frac{\pi}{4}\right]\) 上的最大值为 \(2\)，最小值为 \(-1\).
\end{enumerate}
\end{solution}

\begin{problem}
以下茎叶图记录了甲，乙两组各四名同学的植树棵数. 乙组记录中有一个数据模糊，无法确认，在图中以 \(X\) 表示.

\begin{center}
\begin{tabular}{r|c|l}
甲组 & & 乙组 \\
\hline
\(9\quad 9\) & 0 & \(X\quad 8\quad 9\) \\
\(1\quad 1\) & 1 & \(0\)
\end{tabular}
\end{center}

\begin{enumerate}
\item 如果 \(X = 8\)，求乙组同学植树棵数的平均数和方差；
\item 如果 \(X = 9\)，分别从甲，乙两组中随机选取一名同学，求这两名同学的植树总棵数为 19 的概率.
\end{enumerate}
（注：方差 \( s^2 = \frac{1}{n}\left[(x_1 - \overline{x})^2 + (x_2 - \overline{x})^2 + \ldots + (x_n - \overline{x})^2\right] \)，其中 \(\overline{x}\) 为 \(x_1\)，\(x_2\)，\(\cdots\)，\(x_n\) 的平均数）

\end{problem}

\begin{solution}
\begin{enumerate}
\item 当 \(X=8\) 时，乙组同学植树棵数的平均数为 \(\frac{35}{4}\)，方差为 \(\frac{11}{16}\).
\item 当 \(X=9\) 时，所求概率为 \(\frac{1}{4}\).
\end{enumerate}
\end{solution}

\begin{problem}
如图，在四面体 \(PABC\) 中，\(PC \perp AB\)，\(PA \perp BC\)，点 \(D\)，\(E\)，\(F\)，\(G\) 分别是棱 \(AP\)，\(AC\)，\(BC\)，\(PB\) 的中点.

\begin{enumerate}
\item 求证：\(DE \parallel\) 平面 \(BCP\)；
\item 求证：四边形 \(DEFG\) 为矩形；
\item 是否存在点 \(Q\)，到四面体 \(PABC\) 六条棱的中点的距离相等？说明理由.
\end{enumerate}

\bitmapfigure[max width=0.48\linewidth]{img/2011/beijing_liberal/q17_fig1.png}
\end{problem}

\begin{solution}
\begin{enumerate}
\item 直线 \(DE\) 平行于平面 \(BCP\).
\item 四边形 \(DEFG\) 为矩形.
\item 存在这样的点 \(Q\)，取 \(Q\) 为矩形 \(DEFG\) 的中心即可.
\end{enumerate}
\end{solution}

\begin{problem}
已知函数 \( f(x) = (x - k)\e^x \).
\begin{enumerate}
\item 求 \( f(x) \) 的单调区间；
\item 求 \( f(x) \) 在区间 \( [0,1] \) 上的最小值.
\end{enumerate}
\end{problem}

\begin{solution}
\begin{enumerate}
\item \(f(x)\) 在 \((-\infty,k-1)\) 上单调递减，在 \((k-1,+\infty)\) 上单调递增.
\item 在区间 \([0,1]\) 上的最小值为
\[
\begin{cases}
-k, & k\leqslant 1,\\
-\e^{k-1}, & 1<k<2,\\
(1-k)\e, & k\geqslant 2.
\end{cases}
\]
\end{enumerate}
\end{solution}

\begin{problem}
已知椭圆 \(G: \frac{x^2}{a^{2}} + \frac{y^2}{b^{2}} = 1 \) （\(a > b > 0\)）的离心率为 \(\frac{\sqrt{6}}{3}\)，右焦点为 \((2\sqrt{2}, 0)\)，斜率为 1 的直线 \(l\) 与椭圆 \(G\) 交于 \(A\)，\(B\) 两点. 以 \(AB\) 为底边作等腰三角形，顶点为 \(P(-3, 2)\).
\begin{enumerate}
\item 求椭圆 \(G\) 的方程；
\item 求 \( \triangle PAB \) 的面积.
\end{enumerate}
\end{problem}

\begin{solution}
\begin{enumerate}
\item 椭圆 \(G\) 的方程为 \(\dfrac{x^2}{12}+\dfrac{y^2}{4}=1\).
\item \(\triangle PAB\) 的面积为 \(\dfrac{9}{2}\).
\end{enumerate}
\end{solution}

\begin{problem}
若数列 \(A_{n}\)：\(a_{1}\)，\(a_{2}\)，\(\cdots\)，\(a_{n}\)（\(n \geqslant 2\)）满足 \(\left|a_{k+1} - a_{k}\right| = 1\)（\(k = 1, 2, \cdots, n-1\)），则称 \(A_{n}\) 为 \(E\) 数列. 记 \(S(A_{n}) = a_{1} + a_{2} + \cdots + a_{n}\).
\begin{enumerate}
\item 写出一个 \(E\) 数列 \(A_{n}\) 满足 \(a_{1}=a_{3}=0\)；
\item 若 \(a_1 = 12\)，\(n = 2000\)，证明：\(E\) 数列 \(A_n\) 是递增数列的充要条件是 \(a_n = 2011\)；
\item 在 \(a_1 = 4\) 的 \(E\) 数列 \(A_n\) 中，求使得 \(S(A_n) = 0\) 成立的 \(n\) 的最小值.
\end{enumerate}
\end{problem}

\begin{solution}
\begin{enumerate}
\item 例如 \(A_5=(0,1,0,1,0)\).
\item 充要条件为 \(a_n=2011\).
\item \(n\) 的最小值为 \(9\).
\end{enumerate}
\end{solution}
